%% =====================================================================
\chapter{Administration, Notifications and Settings}\label{ch:admin}

\textbf{Files.} \file{src\codesl{}modules\codesl{}admin\codesl{}admin\codedot{}routes\codedot{}js},
\file{src\codesl{}services\codesl{}notificationService\codedot{}js},
\file{src\codesl{}modules\codesl{}notifications\codesl{}notifications\codedot{}routes\codedot{}js},
\file{src\codesl{}models\codesl{}HospitalSettings\codedot{}js}. Every admin route is behind
\code{router.use(authenticate, requireRole('admin'))}.

\section{Credential verification}

Doctors and insurance agents who register themselves start \emph{pending}.
\code{GET \codesl{}api\codesl{}admin\codesl{}doctors\codesl{}pending-verification} lists them (with licence
number or company), and the admin decides:

\begin{lstlisting}[language=JavaScript,caption={Approving or rejecting (\file{admin.routes.js})}]
  const user = await User.findOne({ _id: req.params.id, role: { $in: User.ROLES_NEEDING_APPROVAL } });
  if (!user) throw notFound('Account not found');
  if (user.approvalStatus !== 'pending') throw conflict(`This account is already ${user.approvalStatus}`, 'ALREADY_DECIDED');

  const approved = req.body.decision === 'approve';
  user.approvalStatus = approved ? 'approved' : 'rejected';
  user.approvalNote = req.body.note;
  await user.save();
  await invalidateSession(user._id); // the new status applies on their next request
  await ledgerService.append({
    chain: 'records', entityType: 'user', entityId: user._id, action: approved ? 'CREDENTIALS_VERIFIED' : 'CREDENTIALS_REJECTED',
    actorId: req.user.userId, actorRole: 'admin', payload: { approvalStatus: user.approvalStatus, licenseNumber: user.doctorProfile?.licenseNumber },
  });
  // ... in-app notification + email to the user
\end{lstlisting}

The decision takes effect on the doctor's very next request, because the
session cache is invalidated (Chapter~\ref{ch:zerotrust}); there is no need for
them to log out. The verification, including which licence number was verified
and by whom, is permanently recorded in the ledger.

\section{User management}

\begin{itemize}
  \item \code{GET \codesl{}api\codesl{}admin\codesl{}users?role=\&status=\&q=} searches by name or email
    and filters by role and active/pending state.
  \item \code{POST \codesl{}api\codesl{}admin\codesl{}users} creates staff accounts (admin, doctor,
    agent) that are approved immediately. This is the only way, besides the
    seed, to create an admin.
  \item \code{PATCH \codesl{}api\codesl{}admin\codesl{}users\codesl{}:id \{ isActive \}} activates or
    deactivates. Deactivation revokes every refresh token and invalidates the
    session cache, so the user is out on their next click. An admin cannot
    deactivate themselves (so the hospital cannot lock itself out by mistake).
\end{itemize}

\section{Live statistics}

\code{GET \codesl{}api\codesl{}admin\codesl{}stats} computes every number on the admin overview at
request time, in one round of parallel queries:

\begin{lstlisting}[language=JavaScript,caption={Parallel aggregation (\file{admin.routes.js}, abridged)}]
  const [usersByRole, pendingVerification, appointmentsToday, appointmentsByStatus, weekly, records, claimsByStatus, revenue,
    devices, readings24h, activeConsents, ledger, mongoPing, redisPing] = await Promise.all([
    User.aggregate([{ $group: { _id: '$role', total: { $sum: 1 }, active: { $sum: { $cond: ['$isActive', 1, 0] } } } }]),
    User.countDocuments({ approvalStatus: 'pending' }),
    // ... appointments today, by status, per day for 7 days
    MedicalRecord.countDocuments(),
    Claim.aggregate([{ $group: { _id: '$status', count: { $sum: 1 }, amount: { $sum: '$amount' } } }]),
    Bill.aggregate([{ $group: { _id: null, billed: { $sum: '$amount' }, paid: { $sum: '$paidAmount' }, covered: { $sum: '$coveredAmount' } } }]),
    // ... devices, readings in 24 h, active consents
    ledgerService.stats(),
    ping(() => mongoose.connection.db.admin().ping()),
    ping(() => redisClient.ping()),
  ]);
\end{lstlisting}

Because nothing is cached or hardcoded, registering a new user changes the user
count on the next refresh; an integration test checks exactly that. The
\code{system} section reports uptime, memory, load and the latency of a MongoDB
and Redis ping, which is useful during a demo to show the infrastructure is live.

\section{Hospital settings}

One settings document (\code{HospitalSettings.getSingleton()}) holds the
hospital name and contacts, the booking window, the default consent duration,
the list of insurers and the departments. The public part is served without
login by \code{GET \codesl{}api\codesl{}settings\codesl{}public}, so the registration form's insurer and
department dropdowns come from the database too. Saving settings is a ledger
entry.

\section{Access audit}

\code{GET \codesl{}api\codesl{}admin\codesl{}audit} reads the \code{access} chain newest first, filtered
by decision (allow/deny) or reason, with the actor's name resolved. It also
returns two counters for the page header: denials in the last 24 hours and the
total number of break-glass accesses. This is the page an auditor uses to answer
``who tried to open what, and were they allowed?''.

\section{IoT device registry}

\code{GET\codesl{}POST\codesl{}PATCH \codesl{}api\codesl{}admin\codesl{}devices} lists devices with their patient and
last-seen time, registers new ones (returning the API key once, Chapter~\ref{ch:rxvitals})
and activates or deactivates them.

\section{Notifications}

\begin{lstlisting}[language=JavaScript,caption={notify and notifyRole (\file{src/services/notificationService.js})}]
/**
 * Creates in-app notifications. Failures are logged and swallowed: a missed
 * notification must never roll back the clinical or financial action that
 * caused it.
 */
async function notify(userIds, { type = 'info', title, message, page = '' }) {
  const ids = [].concat(userIds).filter(Boolean).map(String);
  if (!ids.length) return [];
  try {
    return await Notification.insertMany(ids.map((userId) => ({ userId, type, title, message, page })));
  } catch (error) {
    logger.error(`Notification failed: ${error.message}`);
    return [];
  }
}

async function notifyRole(role, payload) {
  const users = await User.find({ role, isActive: true }).select('_id').lean();
  return notify(users.map((u) => u._id), payload);
}
\end{lstlisting}

\begin{howto}
\code{page} names the dashboard page the notification is about; clicking a
notification in the browser navigates there. \code{notify} accepts one id or a
list, so ``tell every doctor with consent'' is one call. Notifications are
deliberately best-effort, unlike ledger writes: if the ledger write fails, the
request fails; if a notification fails, it is only logged. The dashboards poll
\code{GET \codesl{}api\codesl{}notifications} every 20 seconds to update the bell and sidebar
badges; \code{PATCH \codesl{}:id\codesl{}read} and \code{POST \codesl{}read-all} mark them read.
\end{howto}

\begin{table}[H]
\centering\small
\caption{Who is notified of what}\label{tab:notify}
\begin{tabularx}{\textwidth}{@{}YP{6cm}@{}}
\toprule
\textbf{Event} & \textbf{Recipients} \\
\midrule
Doctor/agent registers & All admins \\
Appointment booked / rescheduled & Doctor \\
Appointment confirmed, completed, cancelled & Patient (or the other party) \\
Access requested & Patient \\
Access granted, approved, denied, revoked & Doctor \\
Break-glass access & Patient and all admins \\
New record, prescription or bill & Patient \\
Refill requested / decided & Doctor / patient \\
Abnormal vital sign & Patient and every doctor with active consent \\
Claim filed & Agents of that insurer \\
Claim status changes & Patient \\
Account verified / rejected & The user (in-app and email) \\
\bottomrule
\end{tabularx}
\end{table}
